\scanpage{346}
% top centre
\mspage{262}
% lines 1–2
\deleted{La condition (i) p. 256 \unsure{implique|indique} que $G(\mathrm{id}_j)=\mathrm{id}_{G(j)}$, la}
% lines 3–7
faire attention que les $t_\beta(\lambda)$ ne définissent \unsure{pas} un hom de systèmes \unsure{projectifs}, mais par la condition (ii) p. 256, ils définissent un hom de $\ast$\unsure{Ind|ind}-objets, donc un hom. sur les $\varinjlim$.
% left margin beside lines 3–10
% slanted
\margin{Attention, c'est ici qu'on utilise \unsure{tacitement} le fait que les foncteurs $\lambda_*:\unsure{\mathcal E|\mathscr E|\mathcal C}_{j'}\longrightarrow\mathcal E_j$ commutent aux $\varinjlim$ de type $I_\alpha$, ce qui implique que les \unsure{inclusions} $\mathcal E_j\longrightarrow\mathcal E$ commutent à ce type de $\varinjlim$.}
% lines 7–12
La condition (i) assure que les $G(\lambda)$ \unsure{transforment|transforme} identités en identités, et la condition (iii) que l'on a $G(\mu\lambda)=G(\mu)G(\lambda)$. Donc $G$ apparaît comme une section de $\mathcal E$ sur $J$. On a d'autre part \deleted{\unsure{des}} des flèches
% line 13
\[
 X(j)\overset{\mathrm{def}}{=}X_0(j)\longrightarrow
 G(j)=\lim_{\beta\in\unsure{I|J}_\alpha}X_\beta(j),
\]
% lines 14–17
ce qui prouve $\unsure{F|P}(\unsure{1|A})$. D'autre part, la condition $3^\circ)$ du lemme prouve la validité de $F'(\unsure{1|A})$. Cela achève la dém. de 4.10.3, et par \unsure{C$^{\mathrm{t}}$|C$'$} de \inserted{4.10.1 et de} 4.9.3.
% archive sheet number: 345
